\documentclass[11pt]{article}
\usepackage[a4paper,margin=26mm]{geometry}
\usepackage{amsmath,amssymb,amsthm,hyperref,booktabs}
\newtheorem{theorem}{Theorem}\newtheorem{lemma}[theorem]{Lemma}
\newcommand{\E}{\mathbb E}\newcommand{\Pp}{\mathbb P}
\title{A constant variance bound removes the distant count tail}
\author{}\date{28 October 2026}
\begin{document}\maketitle

The incoming counts need not be assigned variances growing with their
distance from the own count. Their variances have a common bound. Using
that bound after the first one or two counts produces a factorial tail,
and extends most of the direct comparison to $ns\le .075$. The only
phase interval left out of this note is $-.15<u<0$.

\begin{theorem}\label{thm:outer}
Let $n\ge42$, $2/5\le p\le53/100$, $q=1-p$, and $Y=B+U$, where
$\Pp(B=1)=p$ and $\E(U\mid B)=0$. The two conditional laws satisfy
\[
 |U_j|\le3/5,\quad U_0\le a_0,\quad U_1\ge-a_1,
 \qquad a_j\ge0,\quad a_0+a_1\le9/125.
\]
Put $v=pq$, $s=\E U^2$, $\lambda=ns$, $X=(Y-p)/\sqrt{v+s}$,
$\rho=\E|X|^3$, $\rho_B=(1-2v)/\sqrt v$, and $U_*=3/40$.
Suppose $0<\lambda\le U_*$, $|H|\le46/25$, and $C\ge2/5$.
Choose $k\in\mathbb Z$ with $np+H\sqrt{v+s}=k+u$, where
\[
 u\in[-34/125,-3/20]\ \cup\ [0,91/125].
\]
Let $\delta=k-np$, let $F_B$ be the $\mathrm{Bin}(n,p)$ CDF, and
let $S_n$ be the sum of $n$ independent copies of $X$. Then
\begin{align*}
 \Delta_C={}&\sqrt n\{\Pp(S_n\le H)-\Phi(H/\sqrt n)\}-C\rho\\
 &-\sqrt n\{F_B(k)-\Phi(\delta/\sqrt{nv})\}+C\rho_B
 <-\lambda/100.
\end{align*}
\end{theorem}
The theorem assumes the conditional representation and support bounds;
it does not acquire them. In particular it is not a full Berry--Esseen
theorem, nor does it settle the omitted negative phase. The old $ns\le.03$
comparison covers that omitted phase only at its old noise scale.

\section{Conditioning and normalization}
Write $v_j=\E U_j^2$, and let $W_j$ be the sum of $j$ independent
$U_1$ variables and $n-j$ independent $U_0$ variables. Put
$\pi_j=\binom njp^jq^{n-j}$, $m=\sqrt n\pi_k$, and
\begin{gather*}
 f(u)=\frac{\Phi((\delta+u)/\sqrt{n(v+s)})-
                  \Phi(\delta/\sqrt{n(v+s)})}{\pi_k},\\
 O=\Pp(W_k>u)+f(u),\qquad
 T=\sum_{i=1}^{n-k}\Pp(W_{k+i}\le u-i).
\end{gather*}
For any upper bound $R$ on $\pi_{k+i}/\pi_k$, conditioning gives
\begin{equation}\label{eq:master}
 \Delta_C\le-m(O-RT)+\mathcal R,
 \quad \mathcal R=\sqrt n\{\Phi(\delta/\sqrt{nv})-
 \Phi(\delta/\sqrt{n(v+s)})\}+C(\rho_B-\rho).
\end{equation}
Indeed the own count loses $\pi_k\Pp(W_k>u)$, each larger count
can add only its closed incoming probability, and all other outgoing
terms are nonpositive. This accounts for atoms at every threshold.

Set $A=3/5$, $v_*=6/25$, $N=42$, and
$\Xi=(46/25)(251/500)$. We have $\sqrt{v+s}<251/500$,
$|\delta|<\Xi+|u|\le413/250$, $1\le k\le n-1$, and
$9/25<k/n<23/40$. Finite Stirling expansion gives, for $|\delta|\le D$,
\begin{equation}\label{eq:stirling}
 m=\frac{\phi(0)}{\sqrt v}
       \exp\{-\delta^2/(2nv)+e\},\qquad
 n|e|<Z(D):=5D/8+3/8+D^3/42.
\end{equation}
For completeness, on the segment $[9/25,23/40]$ the third derivative
of the binary relative entropy has absolute value less than $6$,
the derivative of $\log(x(1-x))$ has absolute value less than $5/4$,
and the combined factorial remainder has absolute value less than
$3/(8n)$. Taylor expansion at $p$ gives the three displayed error
terms. The last assertion follows from
$0<\log(m!)-(m+1/2)\log m+m-\tfrac12\log(2\pi)<1/(12m)$
and $x(1-x)>.23$: the absolute difference of the positive remainder
at $n$ and the sum of the two subtracted positive remainders is at
most the larger of their separate upper bounds.
These statements use only $k,n-k>0$. The derivative constants are
checked from $x(1-x)>.23$ and $|1-2x|\le.28$.

For the same fixed $\delta$,
\[
 \log f'(t)=-\tfrac12\log(1+s/v)+\frac{\delta^2}{2nv}
       -\frac{(\delta+t)^2}{2n(v+s)}-e.
\]
It follows, for $|t|\le b$, that
\begin{equation}\label{eq:log}
 n|\log f'(t)|\le\tfrac83\lambda+(2Db+b^2)/(2v_*)+Z(D).
\end{equation}
For a positive phase, the better signed estimate
$2\delta t+t^2\le\max(0,2\Xi u-u^2)$, $0\le t\le u$, may
replace $2Db+b^2$ when proving a lower bound for $f'$.
To see this, write $\delta=\xi-u$, $|\xi|\le\Xi$; the convex
quadratic in $t$ is bounded by its two endpoint values. The resulting
bound increases for $0\le u\le91/125<\Xi$.

Using $.39886<\phi(0)<.4$, $e^{-x}\ge1-x$ and
$e^x\le(1-x)^{-1}$ gives the following rational estimates.
Every table entry is a strict bound.
The value of $f'$ is needed only between $0$ and the actual $u$.
\[
\begin{array}{c|c|c|c|c|c}
 \text{phase}&D&m\text{ lower}&f'&R\text{ upper}&\mathcal R\text{ upper}\\\hline
 {[0,.075]}&1&.73&>.96&1.06&(.75)\lambda/n\\
 {[-.272,-.15]}&1.2&.71&<1.08&1.08&(1.1)\lambda/n\\
 {[.075,.15]}&1.075&.73&>.95&1.07&(.88125)\lambda/n\\
 {[.15,.728]}&1.652&.65&>.9&1.143&(1.891)\lambda/n
\end{array}
\]
Here the binomial ratio at step $j$ equals
$1+(p-\delta-j)/(nv+q(\delta+j))$.
For the first three rows only the first factor can exceed one;
bound it by
$1+(.53+D-1)/(42v_*-.6(D-1))$.
For the last row multiply this bound by
$1+(.53+D-2)/(42v_*)$; all later factors are at most one.

We justify the repair column. Conditional signs are fixed since
$a_0<p$ and $a_1<q$. Cubic expansion and $U_0^3\le a_0U_0^2$,
$U_1^3\ge-a_1U_1^2$ give
\[
 \E|Y-p|^3\ge v(1-2v)+(3\min(p,q)-a_0-a_1)s.
\]
The elementary bound $(1+x)^{3/2}\le1+3x/2+3x^2/8$ for $x\ge0$
therefore yields
\[
 (v+s)^{3/2}(\rho-\rho_B)
 \ge s(87/250-13s/16).
\]
Since $(2/5)(87/250-13U_*/(16N))>(251/500)^3$, we have
$C(\rho-\rho_B)>s$. Differentiating the normal term in its variance
shows it is at most $(7/4)\delta_+s$. Thus
\begin{equation}\label{eq:repair}
 \mathcal R<(7D/4-1)\lambda/n.
\end{equation}
No higher moment of an unconfined law has been used.

\section{A factorial bound for the complete incoming tail}
We use the quadratic-hinge Poisson comparison of Pinelis,
\emph{On the Bennett--Hoeffding inequality},
\href{https://arxiv.org/abs/0902.4058}{arXiv:0902.4058},
Proposition~2.8 and Corollary~2.11. For independent centered summands
bounded above by $A$ and with total variance at most $V$, it gives
\[
 \Pp(Z\ge b)\le
 \frac{\E(P_\mu-j)_+^2}{(b/A+\mu-j)^2},
 \quad \mu=V/A^2,\quad b/A+\mu-j>0,
\]
where $P_\mu$ is Poisson. Variance inflation is allowed by the theorem;
we never assume that the displayed quotient is monotone in $\mu$.

As $|v_1-v_0|\le(5/2)s$, the count at distance $i$ has variance
at most $\lambda[1+(5/2)(D+i)/n]$. More importantly, every legal
count has variance at most
\begin{equation}\label{eq:commonvariance}
 n\max(v_0,v_1)\le(5/2)\lambda.
\end{equation}
Use the first bound for one or two initial counts, then the second
bound for all the rest.

\begin{lemma}\label{lem:tail}
With $\mu=(5/2)\lambda/A^2$, let $m_0\ge2$ be an integer.
If $D_{m_0}:=m_0(1/A-1)-u/A+\mu>0$ and $\mu<m_0+2$, then
\[
 \sum_{i=m_0}^{n-k}\Pp(W_{k+i}\le u-i)
 \le J_{m_0}(\mu,u):=
 \frac{\mu^{m_0+1}}
 {(m_0+1)!D_{m_0}^2(1-\mu/(m_0+2))}.
\]
\end{lemma}
\begin{proof}
Apply the comparison to $-W_{k+i}$, using the common variance bound
and the integer hinge $j=i$. For every integer $l\ge0$ and $i\ge2$,
$(l-i)_+^2\le\binom l{i+1}$. Indeed, writing $l=i+h$, $h\ge1$,
the right side is at least $\binom{h+2}{3}\ge h^2$; the last
inequality is equivalent to $(h-1)(h-2)\ge0$.
Poisson factorial moments then bound the numerator by
$\mu^{i+1}/(i+1)!$. The hinge denominators increase with $i$,
so successive terms have ratio at most $\mu/(i+2)\le\mu/(m_0+2)$.
Summing the infinite geometric majorant proves the claim and includes
all finite incoming counts. Starting this argument at $i=1$ is false:
at $l=3$, $4>\binom32$.
\end{proof}

For the first count we use
\[
 Q(\mu,u)=\frac{\mu^2(1/2+\mu/6)}{((1-u)/A-1+\mu)^2}.
\]
This follows from $\E(P_\mu-1)_+^2=\mu^2-\mu+1-e^{-\mu}
\le\mu^2/2+\mu^3/6$, provided its unsquared denominator is positive.
For the second count, a variance bound $A^2\mu_2$ gives
$\mu_2^3/[6((2-u)/A-2+\mu_2)^2]$.
All these denominators, and all geometric denominators below, are
checked positive on their entire stated domains.

\section{The three outer intervals reduce to endpoint inequalities}
In this section every variance is inflated to its value at $U_*$.
Put $\mu_G=(5/2)U_*/A^2=25/48$. For the negative interval,
Cantelli and the normal derivative bound give
\[
 O\ge L_-:=\min_{r\in\{3/20,34/125\}}
 \left\{\frac{r^2}{(15/14)U_*+r^2}-\frac{27r}{25}\right\}.
\]
The derivative is positive at $r=3/20$ and has only a single maximum
before decreasing, so a minimum on the interval occurs at an endpoint.
The complete tail is at most
\[Q((8/7)U_*/A^2,-3/20)+J_2(\mu_G,-3/20).\]
This uses monotonicity in $u$ only, after fixing the inflated variances.

For the next positive interval $[3/40,3/20]$, use
\[
 O-RT\ge\frac{19u}{20}-\frac{107}{100}
 \left\{Q((9/8)U_*/A^2,u)+J_2(\mu_G,u)\right\}.
\]
The right side is concave in $u$, since its subtracted terms are
positive multiples of reciprocals of squared positive affine functions.
Its two endpoints suffice.

For $3/20\le u\le91/125$, set $D=413/250$,
$t_1=(7/6)U_*$ and $\mu_2=U_*[1+(5/2)(D+2)/42]/A^2$.
Cantelli bounds the first incoming event, giving
\[
 O-RT\ge\frac{9u}{10}-\frac{1143}{1000}
 \left\{\frac{t_1}{t_1+(1-u)^2}
 +\frac{\mu_2^3}{6((2-u)/A-2+\mu_2)^2}+J_3(\mu_G,u)\right\}.
\]
This is again concave: the first tail term has nonnegative second
derivative because $3(1-91/125)^2>t_1$, and the others are convex
reciprocal squares. Thus two more endpoints suffice.
Multiplying by the relevant lower bound for $m$ and subtracting
$(7D/4-1)U_*/42+U_*/100$ gives a strictly positive rational number
in all five tests. Their minimum exceeds $17/10000$.
Each test also proves $O-RT>0$, as needed when replacing $m$ by
its lower bound. The use of $U_*$ in the repair and desired margin
is conservative for every $0<\lambda\le U_*$.

\section{The positive inner interval is one polynomial calculation}
Here $0\le u\le3/40$. Put $\alpha=16/25$, $M=5\alpha^4$ and
$t=\E W_k^2$. The quartic
\[
 h(x)=(-x^4+3\alpha^2x^2+2\alpha^3x)/M
       \le\mathbf1_{(0,\infty)}(x)
\]
has numerator $-x(x+\alpha)^2(x-2\alpha)$ and positive-half-line
maximum $(9+6\sqrt3)/20<1$. Since $\E W_k^3$ has absolute value
at most $At$ and $\E W_k^4\le A^2t+3t^2$, expansion at $W_k-u$
and $f(u)\ge(24/25)u$ yield
\begin{equation}\label{eq:quartic}
 O\ge L(t,u):=
 \frac{(3\alpha^2-A^2)t-3t^2-4Aut-6u^2t-u^4
       +3\alpha^2u^2-2\alpha^3u}{M}+\frac{24u}{25}.
\end{equation}
This is concave in $t$, and
$(79/84)\lambda\le t\le(89/84)\lambda$.
Use the actual $\lambda$ in the tail, with
$\mu_1=25\lambda/8$ and $\mu_G=125\lambda/18$. It remains to show,
for each $c\in\{79/84,89/84\}$,
\begin{equation}\label{eq:innergoal}
 \frac{73}{100}\{L(c\lambda,u)-\tfrac{53}{50}
 [Q(\mu_1,u)+J_2(\mu_G,u)]\}
       -(\tfrac{3}{168}+\tfrac1{100})\lambda>0.
\end{equation}
Here the $3/168$ is the normal repair after the unit cubic gain.

This last verification is a small exact polynomial certificate. Set
\[
 d_1=2/3-5u/3+25\lambda/8,\quad
 d_2=4/3-5u/3+125\lambda/18,\quad d_3=1-125\lambda/72.
\]
Multiply the left side of \eqref{eq:innergoal} by
$6d_1^2d_2^2d_3>0$. The result $P_c(\lambda,u)$ has degree at
most $7$ in $\lambda$ and $8$ in $u$. Write
$P_c=\sum a_{kl}\lambda^ku^l$. Its tensor Bernstein coefficients
on the single square $[0,3/40]^2$ are exactly
\[
 b_{ij}=\sum_{k\le i,l\le j}a_{kl}(3/40)^{k+l}
       \frac{\binom ik}{\binom7k}\frac{\binom jl}{\binom8l},
 \qquad 0\le i\le7,\quad0\le j\le8.
\]
The source \texttt{outer-check.py} constructs these polynomials from
\eqref{eq:quartic}--\eqref{eq:innergoal} using rational arithmetic,
and checks all $144$ coefficients: $b_{00}=0$ and every other
coefficient is strictly positive, for both values of $c$.
The transform follows from
$x^k=\sum_{i=k}^m\binom ik/\binom mk\,B_i^m(x)$.
Bernstein basis functions are nonnegative and sum to one. At every
point of the square except $(0,0)$ at least one nonzero-index basis
function is positive. Consequently $P_c>0$ whenever $\lambda>0$,
including all boundary phases. This proves \eqref{eq:innergoal} and
completes the theorem.

\section{What the check does and does not establish}
The acceptance script checks the scalar normalization and endpoint
inequalities and reconstructs both coefficient arrays. Missing, repeated,
reordered and corrupted arrays are rejected; removing the cubic gain
or increasing the noise endpoint to $.1$ fails a genuine sufficient
inequality. The proof above supplies the analytic reduction; an arithmetic
PASS alone supplies neither that reduction nor acquisition of its premises.
The independent review and saved runs are separate files. No inherited
certificate or manuscript is replaced by this component.
\end{document}
